\section{Background}

In this section we introduce the background necessary to understand the contributions of this paper. 

\subsection{\mk}

The \mk family of programming languages is created to be simple to understand, implement, and extend with new features.\footnote{Website of the \mk programming languages family: \url{http://minikanren.org/}} 
One distinguishing characteristic of \mk is its complete search enabled by design. 
Any existing solution to a query will be found eventually~\cite{rozplokhas2020certified} which makes \mk uniquely suitable for declarative setting and reverse execution. 

Most \mk dialects are implemented as a set of combinators for basic operations such as a unification or conjunction and disjunction. 
However, program transformation is complicated for shallow embeddings in languages featuring referential transparency. 
Primarily due to this fact the proof-of-concept functional conversion was implemented for a deeply embedded \mk and manipulated abstract syntax trees explicitly. 
The syntax in Fig.~\ref{fig:minikanren} describes the core language used in this paper, also known as \micro~\cite{friedman2013mukanren}. 
We also assume inverse $\eta$-delay to accompany every relation invocation and not used elsewhere.  

\begin{figure}[h]
\begin{tabular}{llll}
    $\mathcal{C}$ & $=$ & $\{\Cons_{i}^{k_{i}}\} $ & constructors with arities \\
    $\mathcal{R}$ & $=$ & $\{R_{i}^{k_{i}}\} $ & relational symbols with arities \\
    $\Term_{V}$ & $=$ & $V \cup \{\Cons_{i}\left( t_1, \ldots, t_{k_{i}} \right) \mid t_{j}\in \Term_{V}\} $ & terms on set of variables $V$ \\
 $\Kan_{V}$ & $=$ & $\Term_{V} \equiv \Term_{V}$ & unification \\
 & $\mid$ & $\Kan_{V} \land \Kan_{V}$ & conjunction  \\
 & $\mid$ & $\Kan_{V} \lor \Kan_{V}$ & disjunction  \\
 & $\mid$ & $R_{i}^{k_{i}}\left( t_1, \ldots, t_{k} \right), t_{j}\in \Term_{V}$ & relational symbol invocation\\
 & $\mid$ & $\Fresh_{x} \Kan_{V}$ & fresh variable introduction  \\
 $\Def_{V}$ & $=$ & $R_{i}\left( x_1, \ldots, x_{k_{i}} \right) \equiv \Kan_{V}, x_{j}\in V$ & relational symbol definition \\
  $\mathcal{P}$ & $=$ & $\Def_{V}^*$ & miniKanren program
\end{tabular}
\caption{miniKanren syntax}
\label{fig:minikanren}
\end{figure}

The addition relation written in this syntax takes the following shape, where $R_{add}$ is the only relational symbol, $\Cons_{0}$ and $\Cons_{S}$ are constructors, and $x, x', y, z, z'$ are variables:

\begin{equation*}
    \begin{split}
        R_{\text{add}}\left( x, y, z \right) \equiv 
        &\left(  x \equiv \Cons_0  \land y \equiv z \right) \\
        \lor &\left( \Fresh_{x'}\Fresh_{z'} \left(x \equiv \Cons_{S}\left( x' \right)  \land R_{add}\left(x', y, z'\right) \land z \equiv \Cons_{S}\left( z' \right) \right)\right)
    \end{split}
\end{equation*}

\subsection{Normal Form}
\label{background:normal}

Many program transformations become simpler when done over a normalized representation, such as the superhomogenous normal form used in the Mercury programming language~\cite{mercuryexec}. 
It is in essence a disjunctive normal form with unique variables, and with extra care taken not to allow exponential explosion in the number of clauses. 
To avoid it, we create  a new relation whenever a disjunction takes place within a conjunction. 
The process of normalization is described in detail in \cite{engel2023translation}, here we only present the normal form syntax, see Figure \ref{fig:normalizedminikanren}.
Notice that there are also no $\Fresh$ declarations in this form, and each element of $V$ is assumed to refer to the same variable within the invocation. 

\begin{figure}[h]
\begin{tabular}{llll}
$\FlatTerm_{V}$ & $=$ & $V \cup \{\Cons_{i}\left( x_1, \ldots, x_{k_{i}} \right) \mid x_{j}\in V, a \neq b \implies x_{a} \neq x_{b}\}$ & linear flat terms\\
$\Base_{V}$ & $=$ & $V \equiv \FlatTerm_{V}$ & linear flat unification\\
            & $\mid$ & $R_{i}\left( x_1, \ldots, x_{k_{i}} \right), x_{j}\in V, a \neq b \implies x_{b} \neq x_{b}$ & linear flat call\\
$\Conj_{V}$ & $=$ & $\bigwedge\left( g_1, \ldots, g_n \right), g_{i}\in \Base_{V}$ & normal conjunction \\
$\KanN_{V}$ & $=$ & $\bigvee\left( c_1, \ldots, c_{n} \right), c_{i}\in \Conj_{V}$ & normal form
\end{tabular}
\caption{Normalized miniKanren syntax}
\label{fig:normalizedminikanren}
\end{figure}


\subsection{Modes}

A relational program can be executed in multiple directions, also called \emph{modes}. 
Mode systems assign each variable in a relation a mode annotation that denotes how the variable is used in it. 
A logic/functional programming language \merc features an advanced mode system that informs its compiler on possible optimizations. 
The functional conversion of \mk relies on the simplest mode system with only two possible annotations: $\In$ and $\Out$. 
$\In$ denotes a variable whose value is fully known when the operation is executed. 
$\Out$ is reserved for a variable with unknown value which is going to be bound after the operation is executed. 

When all variables of a relation are annotated with modes, we say that the relation has a mode. 
We use the term \emph{direction} to describe a relation call whose arguments are annotated. 
For example, consider the following directions of the addition relation:

\begin{equation*}
    \begin{split}
        R_{\text{add}}\left( x^{\In}, y^{\In}, z^{\Out} \right) \equiv 
        &\left(  x^{\In} \equiv \Cons_0  \land y^{\In} \equiv z^{\Out} \right) \\
\lor \left( \Fresh_{x'}\Fresh_{z'} \left(x^{\In} \equiv \Cons_{S}\left( x'^{\Out} \right)\right.\right.  \land& \left.\left. R_{add}\left(x'^{\In}, y^{\In}, z'^{\Out}\right) \land z^{\Out} \equiv \Cons_{S}\left( z'^{\In} \right) \right)\right)
    \end{split}
\end{equation*}
\begin{equation*}
    \begin{split}
        R_{\text{add}}\left( x^{\Out}, y^{\Out}, z^{\In} \right) \equiv 
        &\left(  x^{\Out} \equiv \Cons_0  \land y^{\Out} \equiv z^{\In} \right) \\
\lor \left( \Fresh_{x'}\Fresh_{z'} \left(x^{\Out} \equiv \Cons_{S}\left( x'^{\In} \right)\right.\right.  \land& \left.\left. R_{add}\left(x'^{\Out}, y^{\Out}, z'^{\In}\right) \land z^{\In} \equiv \Cons_{S}\left( z'^{\Out} \right) \right)\right)
    \end{split}
\end{equation*}
\begin{equation*}
    \begin{split}
        R_{\text{add}}\left( x^{\In}, y^{\Out}, z^{\Out} \right) \equiv 
        &\left(  x^{\In} \equiv \Cons_0  \land y^{\Out} \equiv z^{\Out} \right) \\
\lor \left( \Fresh_{x'}\Fresh_{z'} \left(x^{\In} \equiv \Cons_{S}\left( x'^{\Out} \right)\right.\right.  \land& \left.\left. R_{add}\left(x'^{\In}, y^{\Out}, z'^{\Out}\right) \land z^{\Out} \equiv \Cons_{S}\left( z'^{\In} \right) \right)\right)
    \end{split}
\end{equation*}

In a well-moded relation, data flows from $\In$ to $\Out$, and once a variable's value is known, it is never forgotten.
This means that a variable with the $\In$ mode can never become $\Out$, while $\Out$ becomes $\In$ after the current operation is executed.
This intuition translates well to functional code, where each variable is expected to be bound by some value before use.
The functional conversion employs a mode analysis capable of inferring mode annotations given a direction.
The mode analysis algorithm does not fall in the scope of this paper; thus, the reader is referred to~\cite{engel2023translation} for a detailed description of the approach.
Its details are not essential for understanding the contributions presented here.

\subsection{Functional Conversion}

In this section we briefly summarize the functional conversion described in~\cite{engel2023translation} which the reader can refer to for deeper comprehension. 
The functional conversion mirrors the relational conversion~\cite{lozov2018typed} aimed at generating relations from functions thus coming full circle in the verifier-to-solver approach.

Given a \mk relation with a concrete direction, the goal of a functional conversion is to construct a function which yields the same answers as the relation would. 
Because the search in \mk is complete, disjunctions and conjunctions can be reordered, which will hopefully result in faster execution. 
Mode analysis guides this reordering as well as classifies unifications into equality checks, assignments, pattern matches and generations. 
Consider the following functional counterpart of $R_{\text{add}}\left( x^{\In}, y^{\Out}, z^{\Out} \right)$ which computes such pairs $(y, z)$ that $x + y = z$, where $x$ is known. 

\includelines{code/addoIOO.hs}{0}{18}

The function \lstinline{addoIOO} produces infinitely many pairs of Peano numbers; thus the output of the function should be a stream of \lstinline{(Term, Term)}. 
The \lstinline{Stream} is not always the most efficient data structure to represent the output, especially when the computation is deterministic. 
Because of this, the converter generates functions parameterized by an arbitrary \lstinline{Monad m} that can be specified at the call site. 

Let us now illustrate the conversion process by this example. 
Once mode analysis is finished, disjunctions are translated into \lstinline{msum} of several monadic computations (line 6). 
Each such computation is produced from a conjunction of either calls or unifications. 
Note that moded unifications can become an equality test (line 8), a pattern matching (line 13), an assignment (line 15), or use term generation (line 9). 
Generation is required whenever both sides of a unification are annotated $\Out$, producing a stream of all possible values. 

A drawback of the previous implementation is the need to supply generators as explicit function arguments. 
Moreover, when several variables are generated, a separate generator must be provided for each one—even if they share the same type and could be produced in exactly the same way. 
This requirement makes the resulting code inelegant, as well as prone to user errors when they later have to supply the generators manually.




